\chapter{Building and remaking a set}\label{ch:sets}

\section{Step by step in the app}
\begin{enumerate}
\item \textbf{Library}: press \emph{Add audio files} and choose your tracks, or drop them on the panel. Wait for each card to say \emph{ready}.
\item Read the \textbf{pair table} (Figure~\ref{fig:matrix}): green pairs blend easily, amber with care, red are better cut or echoed out.
\item Tick \emph{in the set} on each track you want.
\item \textbf{Set builder}: choose a \emph{Transition style}, or leave \emph{Auto}; press \emph{Auto order}, or move tracks with \emph{Up}.
\item Read each transition (Figure~\ref{fig:set}): where it leaves the old track, the speed and pitch change, the key after it, and why the style was chosen. Open \emph{The rule behind this preset} for the scene's habit.
\item Press \emph{Render MP3} (or WAV). The app plays the whole set through its filters faster than real time, encodes it, and saves the audio and a set file.
\end{enumerate}

\begin{figure}[H]\centering
\includegraphics[width=\linewidth]{matrix.png}
\caption{How every pair of tracks would mix: speed change for the second track and key fit after it.}\label{fig:matrix}
\end{figure}

\begin{figure}[H]\centering
\includegraphics[width=\linewidth]{set.png}
\caption{The Set builder with the four tracks in automatic order, each transition explained.}\label{fig:set}
\end{figure}

\section{Set 1: Alfaaz and Faded}
\begin{tabularx}{\linewidth}{@{}Xrlrr@{}}\toprule
Track & BPM & Key & RMS (dB) & Length \\\midrule
Hamza Malik, Zain Zohaib - Alfaaz & 80.1 & A major (11B) & -13.5 & 3:04 \\
Alan Walker - Faded & 89.9 & F\# major (2B) & -8.5 & 3:33 \\
\bottomrule\end{tabularx}

\begin{enumerate}
\item \textbf{Hamza Malik, Zain Zohaib - Alfaaz $\to$ Alan Walker - Faded}: Any style: echo out. Leave at 2:48; next track at 1$\times$ (0\,\%), pitch 0 semitones, key after the change 2B: 3 steps apart: melodies will clash, so mix on drums or after the vocals end. \emph{Why:} tempos 80 and 90 BPM are 11\% apart even at half or double time.
\end{enumerate}

Tempos 80 and 90\,BPM are 11\,\% apart, and even at half or double time no match is closer, so blending would audibly bend both voices. The plan echoes \emph{Alfaaz} out on a phrase boundary and starts \emph{Faded} clean on its first beat.

\section{Set 2: Socha Hai and It's the Time to Disco}
\begin{tabularx}{\linewidth}{@{}Xrlrr@{}}\toprule
Track & BPM & Key & RMS (dB) & Length \\\midrule
Farhan Akhtar - Socha Hai & 134.2 & D major (10B) & -10.5 & 4:06 \\
Vasubdhara Das Kay Kay Shaan Loy Mendonsa - It's the time to Disco & 132.5 & F\# minor (11A) & -10.9 & 5:35 \\
\bottomrule\end{tabularx}

\begin{enumerate}
\item \textbf{Farhan Akhtar - Socha Hai $\to$ Vasubdhara Das Kay Kay Shaan Loy Mendonsa - It's the time to Disco}: House and EDM: long blend with a bass swap. Leave at 3:20; next track at 1.013$\times$ (1.3\,\%), pitch 0.22 semitones, key after the change 11A: a diagonal move (one step and a change between major and minor): most notes shared, a change of mood. \emph{Why:} tempos within 1.3\% and the keys fit, so a long blend works.
\end{enumerate}

Here the tempos are within about 1.3\,\%, a change of a fifth of a semitone, and the keys sit diagonally on the wheel, so a long blend works. Both are Bollywood songs: if the blend sounds busy because both have vocals, choose the \emph{Bollywood and desi} style instead, which cuts cleanly on the sam.

\section{Both sets together}
\begin{tabularx}{\linewidth}{@{}Xrlrr@{}}\toprule
Track & BPM & Key & RMS (dB) & Length \\\midrule
Hamza Malik, Zain Zohaib - Alfaaz & 80.1 & A major (11B) & -13.5 & 3:04 \\
Alan Walker - Faded & 89.9 & F\# major (2B) & -8.5 & 3:33 \\
Vasubdhara Das Kay Kay Shaan Loy Mendonsa - It's the time to Disco & 132.5 & F\# minor (11A) & -10.9 & 5:35 \\
Farhan Akhtar - Socha Hai & 134.2 & D major (10B) & -10.5 & 4:06 \\
\bottomrule\end{tabularx}

\begin{enumerate}
\item \textbf{Hamza Malik, Zain Zohaib - Alfaaz $\to$ Alan Walker - Faded}: Any style: echo out. Leave at 2:48; next track at 1$\times$ (0\,\%), pitch 0 semitones, key after the change 2B: 3 steps apart: melodies will clash, so mix on drums or after the vocals end. \emph{Why:} tempos 80 and 90 BPM are 11\% apart even at half or double time.
\item \textbf{Alan Walker - Faded $\to$ Vasubdhara Das Kay Kay Shaan Loy Mendonsa - It's the time to Disco}: Any style: echo out. Leave at 2:51; next track at 1$\times$ (0\,\%), pitch 0 semitones, key after the change 11A: 3 steps apart: melodies will clash, so mix on drums or after the vocals end. \emph{Why:} tempos 90 and 133 BPM are 32\% apart even at half or double time.
\item \textbf{Vasubdhara Das Kay Kay Shaan Loy Mendonsa - It's the time to Disco $\to$ Farhan Akhtar - Socha Hai}: House and EDM: long blend with a bass swap. Leave at 5:04; next track at 0.9872$\times$ (-1.28\,\%), pitch -0.22 semitones, key after the change 10B: a diagonal move (one step and a change between major and minor): most notes shared, a change of mood. \emph{Why:} tempos within 1.3\% and the keys fit, so a long blend works.
\end{enumerate}

The automatic order starts with the calmest track and moves through the slower pair into the faster one. The first two changes are echo-outs because the gap from 90 to 133\,BPM cannot be bridged without changing the songs; between the two fast tracks it blends.

\section{The set file and remaking a mix}
Rendering also saves \texttt{mix\_DATE\_set.json}. It records the app version, every track's name, length, measurements and \textbf{SHA-256 fingerprint} (a 64-character code that changes completely if one byte of the file changes), every transition with its speeds, times and automation points, and the timeline: when each track starts, from where, at what speed and with what loudness trim. To remake the mix, add the same files to the Library, press \emph{Open set file}, and render. The audio itself is never stored in the set file, so a set can be shared without sharing the music.

From the command line, the same planning runs without a browser:
\begin{lstlisting}[language=bash]
node tools/plan.mjs --order auto file1.mp3 file2.mp3 file3.mp3
node tools/plan.mjs --preset bollywood-cut "Socha Hai.mp3" "Disco.mp3"
\end{lstlisting}
